\documentclass[12pt,aspectratio=169]{beamer}
\input{header_beamer}
\usetheme{Boadilla}
\usecolortheme{beaver}
\setbeamertemplate{page number in head/foot}{}
\setbeamertemplate{navigation symbols}{}
\title{Lecture-24: Poisson Point Processes}
\date{}

\begin{document}
\pagestyle{empty}
\maketitle



\begin{frame}{Simple point processes} 
Consider the $d$-dimensional Euclidean space $\R^d$. 
The collection of Borel measurable subsets $\sB(\R^d)$ of the above Euclidean space is generated by sets $B(x) \triangleq \set{y \in \R^d: y_i \le x_i}$ for $x \in \R^d$.  
\begin{Definition}<2->
A \textbf{simple point process} is a random countable collection of \emph{distinct points} $S: \Omega \to \sX^\N$, 
such that the distance $\norm{S_n} \to \infty$ as $n \to \infty$. 
\end{Definition}
\begin{block}{Reminder}<3->
Since $S$ is a simple point process, each point $S_n$ is unique. 
Therefore, we can identify $S$ as a random set of points in $\sX$ and $S\cap A$ is the random set of points in $A$. 
\end{block}
\end{frame}

\begin{frame}
\begin{block}{Reminder}
For any simple point process $S$, we have $P(\set{S_n = S_m \text{ for any } n \neq m}) = 0$ and $\abs{S\cap A}$ is finite almost surely for any bounded set $A \in \sB(\sX)$.
\end{block}

\begin{Example}[Simple point process on the half-line]<2->
We can simplify this definition for $d=1$. 
When $\sX = \R_+$, one can order the points of the process $S: \Omega \to \R_+^\N$ to get ordered process $\tilde{S}: \Omega \to \R_+^\N$, 
such that $\tilde{S}_n = S_{(n)}$ is the $n$th order statistics of $S$. 
That is, $S_{(0)} \triangleq 0$, 
and 
$
S_{(n)} \triangleq \inf\set{S_k > S_{(n-1)}: k \in \N}. 
$ 
such that $S_{(1)} < S_{(2)} < \dots < S_{(n)} < \dots$, and $\lim_{n \in \N}S_{(n)} = \infty$. 
We will call this an arrival process. 
\end{Example}
\end{frame}

\begin{frame}
\begin{Definition}
Corresponding to a point process $S:\Omega\to\sX^\N$, 
we denote the number of points in a set $A \in \sB(\sX)$ by
\EQ{
N(A)\triangleq \abs{S\cap A} = \sum_{n \in \N}\Ind{A}(S_n), \text{ where we have }N(\emptyset) =0. 
}
The resulting process $N: \Omega \to {\Z_+}^ {\sB(\sX)}$ is called a \textbf{counting process} for the point process $S: \Omega \to \sX^\N$. 
\end{Definition} 
\begin{block}{Reminder}<2->
Let $A \in \sB(\sX)^k$ be a bounded partition of $B\in \sB(\sX)$. 
From the disjointness of $(A_1, \dots, A_k)$, we have 
\EQ{
N(B) = \sum_{n \in \N}\Ind{\cup_{i=1}^kA_i}(S_n)
= \sum_{n \in \N}\sum_{i=1}^k\Ind{A_i}(S_n)
= \sum_{i=1}^k\sum_{n \in\N}\Ind{A_i}(S_n)
= \sum_{i=1}^kN(A_i).
} 
\end{block}
\end{frame}

\begin{frame}
\begin{Definition} 
A counting process is \textbf{simple} if the underlying point process is simple. %, and $N(\set{x}) \in \set{0,1}$ for all $x \in \R^d$. 
\end{Definition}
\begin{block}{Reminder}<2->
For a simple counting process $N$, we have $N(\set{x}) \le 1$ almost surely for all $x \in \sX$.
\end{block}
\begin{block}{Reminder}<3->
Let $N : \Omega \to {\Z_+}^{\sB(\sX)}$ be the counting process for the point process $S: \Omega \to \sX^\N$.  
\begin{enumerate}[i\_]
\item Note that the point process $S$ and the counting process $N$ carry the same information. 
\item The distribution of point process $S$ is completely characterized by the finite dimensional distributions of random vectors $(N(A_1), \dots, N(A_k))$ for any bounded sets $A_1, \dots, A_k \in \sB(\sX)$ and finite $k \in \N$.  
\end{enumerate}
\end{block}
\end{frame}

\begin{frame}
\begin{Example}[Simple point process on the half-line] 
Since the Borel measurable sets $\sB(\R_+)$ are generated by half-open intervals $\set{(0,t]: t \in \R_+}$, 
we denote the counting process by $N: \Omega \to {\Z_+}^{\R_+}$, 
where $N_t \triangleq N(0, t]= \sum_{n \in \N}\SetIn{S_n \in (0,t]}$ is the number of points in the half-open interval $(0,t]$. 
For $s < t$, the number of points in interval $(s,t]$ is $N(s,t] = N(0,t] - N(0,s] = N_t - N_s$. 
\end{Example}

\begin{block}{Theorem : R\'{e}nyi}<2->
\label{thm:void}
Distribution of a simple point process $S:\Omega\to\sX^\N$ on a locally compact second countable space $\sX$ is completely determined by void probabilities $(P\set{N(A)=0}: A \in \sB(\sX))$. 
\end{block}
\end{frame}


\begin{frame}
\begin{block}{Proof}
It suffices to show that the finite dimensional distributions of $S$ on locally compact sets are characterized by void probabilities. 
\begin{enumerate}
\item[\bf{Step 1:}]  
We will show this by induction on the number of points $k$ in a bounded set $A \in \sB$. 
Let $A_1, \dots, A_k, B \in \sB(\sX)$ locally compact, then we will show that $u_k \triangleq P(\cap_{i=1}^k\set{N(A_i)> 0}\cap\set{N(B) = 0})$ can be computed from void probabilities.
From $k=1$, we have 
\EQ{
P\set{N(A_1) > 0, N(B) = 0} = P\set{N(B) = 0} - P\set{N(B\cup A_1) = 0}.
}
The induction can be proved by the recursive relation 
\begin{align*}
u_k 
&= P(\cap_{i=1}^{k-1}\set{N(A_i)> 0}\cap\set{N(B) = 0}) \\
&- P(\cap_{i=1}^{k-1}\set{N(A_i)> 0}\cap\set{N(A_k\cup B) = 0}).
\end{align*}
\end{enumerate}
\end{block}
\end{frame}

\begin{frame}
\begin{block}{Proof}
\begin{enumerate}
\item[\bf{Step 2:}]  
For any locally compact set $B \in \sB(\sX)$, there exists a sequence of nested partitions $B_n \triangleq (B_{n,j}: j \in [J_n])$ that eventually separates the points in $S\cap B$ as $n \to \infty$. 
We define the number of subsets of partition $(B_{n,j}: j \in [J_n])$ that consist of at least one point in $S\cap B$, as $H_n(B) \triangleq \sum_{j=1}^{J_n}\SetIn{N(B_{n,j}) > 0}$ where $H_n(B) \uparrow N(B)$ almost surely. 
\end{enumerate}
\end{block}
\end{frame}

\begin{frame}
\begin{Proof}
\begin{enumerate}
\item[\bf{Step 3:}]  
We next show that for all locally compact sets $B_1, \dots, B_k \in \sB(\sX)$ and $j_1, \dots, j_k \in \N$, 
the probability $P(\cap_{i=1}^k\set{H_n(B_i) = j_i})$ can be expressed in terms of void probabilities. 
We observe that 
\eq{
P(\cap_{i=1}^k\set{H_n(B) = j_i})\\
= \sum_{ T_1, \dots, T_k \subseteq [J_n]: \abs{T_1} = j_1, \dots, \abs{T_k}=j_k}&P\Big(\cap_{i=1}^k\cap_{j\in T_i}\set{N(B^i_{n,j})>  0}\cap \\
&\set{N(\cup_{j \notin \cup_{i=1}^kT_i}B^i_{n,j}) =0}\Big).
}
This can be expressed in terms of void probabilities by Step~1. 
\item[\bf{Step 4:}]<2-> For a simple point process, we have the following almost sure limit $\lim_n\cap_{i=1}^k\set{H_n(B_i) = j_i} = \cap_{i=1}^k\set{N(B_i) = j_i}$. 
The result follows from the continuity of probability.
\end{enumerate}
\end{Proof}
\end{frame}


\begin{frame}
\begin{block}{Reminder} 
Recall that $\abs{A} = \int_{x \in A}dx$ is the volume of the set $A \in \sB(\R^d)$ and for any such $A$.
\end{block}
\begin{Definition}<2->
The \textbf{intensity measure} $\Lambda:\sB(\sX)\to\R_+$ is defined for each bounded set $A \in \sX$ as its scaled volume 
in terms of the \textbf{intensity density} $\lambda: \R^d \to \R_+$, as 
\EQ{
\Lambda(A) \triangleq \int_{x \in A}\lambda(x) dx. 
} 
If the intensity density $\lambda(x) = \lambda$ for all $x \in \R^d$, then $\Lambda(A) = \lambda\abs{A}$. 
In particular for partition $A_1, \dots, A_k$ for a set $B$, we have $\Lambda(B) = \sum_{i=1}^k\Lambda(A_i)$. 
\end{Definition}
\end{frame}


\begin{frame}{Poisson point process}
\begin{Definition}
A non-negative integer valued random variable $N: \Omega \to \Z_+$ is called \textbf{Poisson} if 
for some constant $\lambda > 0$, we have 
\EQ{
P\set{N = n} = e^{-\lambda}\frac{\lambda^n}{n!}.
}
\end{Definition}
\begin{block}{Reminder}<2->
It is easy to check that $\E N = \Var[N] = \lambda$. 
Furthermore, the moment generating function $M_{N_t} = \E e^{t N} = e^{\lambda(e^{t} - 1)}$ exists for all $t \in \R$. 
\end{block}
\end{frame}


\begin{frame}
\begin{cor}
\label{cor:ExpPoisson}
A simple counting process $N: \Omega \to \Z_+^{\sB(\sX)}$ has Poisson marginal distribution with intensity measure $\Lambda:\sB(\sX)\to\R_+$ if and only if void probabilities are exponential with the same intensity measure $\Lambda$. 
\end{cor}
\begin{block}{Proof}<2->
It is clear that if the marginal distribution of the counting process $N$ is Poisson with intensity measure $\Lambda$, then the void probability $P\set{N(A) = 0} = e^{-\Lambda(A)}$ is exponential for any bounded set $A \in \sB(\sX)$. 
\end{block}
\end{frame}


\begin{frame}
\begin{block}{Proof continued}
Conversely, we assume that the void probabilities are exponentially distributed with intensity measure $\Lambda$. 
It follows from the linearity of intensity measure that for any finite, bounded, and disjoint sets $B_1, \dots, B_k \in \sB(\sX)$, we have 
\begin{align*}
P(\cap_{i=1}^k\set{N(B_i)=0}) 
&= P\set{N(\cup_{i=1}^kB_i) = 0}
= e^{-\Lambda(\cup_{i=1}^kB_i)} \\
&= \prod_{i=1}^ke^{-\Lambda(B_i)}
= \prod_{i=1}^kP\set{N(B_i) = 0}.
\end{align*}
\end{block}
\end{frame}


\begin{frame}
\begin{block}{Proof continued}
\begin{itemize}
\item That is, the Bernoulli random vector $(\SetIn{N(B_i) = 0}: i \in [k])$ is independent for any finite $k \in \N$ and bounded disjoint $\sB(\sX)$ measurable sets $B_1, \dots, B_k$. 
Next we consider a set $B \in \sB(\sX)$ and a partition $B_n \triangleq (B_{n,j}: j \in [J_n])$ of $B$ such that $\Lambda(B_{n,j}) = \frac{\Lambda(B)}{J_n}$ for all $j \in [J_n]$. 
\item <2-> It follows that $H_n(B) \triangleq \sum_{j=1}^{J_n}\SetIn{N(B_{n,j})> 0}$ is the sum of $J_n$ \iid Bernoulli random variables with success probability $p_n \triangleq 1-e^{-\Lambda(B)/J_n}$, 
and hence has a Binomial distribution with parameters $(J_n, p_n)$.  
\end{itemize}
\end{block}
\end{frame}

\begin{frame}
\begin{Proof}[Proof continued]
\begin{itemize}
\item Therefore, 
\begin{align*}
P\set{H_n(B) = m} 
&= \frac{e^{-\Lambda(B)}}{m!}\binom{J_n}{m}(e^{\Lambda(B)/J_n}p_n)^m \\
&= e^{-\Lambda(B)}\frac{J_n!}{(J_n-m)!}\Big(e^{\Lambda(B)/J_n}-1\Big)^m.
\end{align*}
\item <2-> Recall that $H_n(B) \uparrow N(B)$ as $n\to\infty$ in the proof of R\'{e}nyi's Theorem, 
and $\lim_{n\to\infty}J_n = \infty$ and $\lim_{n\in\N}\abs{B_{n,j}}= 0$. 
Thus, $\lim_{n\to\infty}\frac{J_n!}{(J_n-m)!}(e^{\Lambda(B)/J_n}-1)^m = \Lambda(B)^m$.
Taking limit $n\to\infty$ on both sides of the above equation, we get the result. 
\end{itemize}
\end{Proof}
\end{frame}

\begin{frame}
\begin{Definition} 
A counting process $N:\Omega\to\Z_+^{\sB(\sX)}$ has the \textbf{completely independence property},  
if for any collection of finite disjoint and bounded sets $A_1, \dots, A_k \in \sB(\sX)$, 
the vector $(N(A_1), \dots, N(A_k)):\Omega \to \Z_+^k$ is independent. 
That is,
\EQ{
P\left(\bigcap_{i=1}^k\set{N(A_i) = n_i}\right) = \prod_{i=1}^kP\set{N(A_i) = n_i},\quad n \in \Z_+^k.
}
\end{Definition}

\begin{Definition}<2->
A simple point process $S: \Omega \to \sX^\N$ is \textbf{Poisson point process}, 
if the associated counting process $N: \Omega \to \Z_+^{\sB(\sX)}$ has complete independence property and the marginal distributions are Poisson. 
\end{Definition}
\end{frame}


\begin{frame}
\begin{Definition}
The \textbf{intensity measure} $\Lambda: \sB(\sX) \to \R_+$ of Poisson process $S$ is defined by $\Lambda(A) \triangleq \E N(A)$ for all bounded $A \in \sB(\sX)$. 
\end{Definition}

\begin{block}{Reminder}<2->
Recall that for any partition $A \in \sB(\sX)^k$ of a bounded set $B \in \sB(\sX)$, we have $N(B) = \sum_{i=1}^kN(A_i)$ and therefore it follows from the linearity of expectations that $\Lambda(B) = \E N(B) = \sum_{i=1}^k\E N(A_i) = \sum_{i=1}^k\Lambda(A_i).$ 
Thus, this is a valid intensity measure. 
\end{block}
\end{frame}

\begin{frame}
\begin{block}{Reminder}
For a Poisson process with intensity measure $\Lambda$, 
it follows from the definition that for any finite $k \in \Z_+$, 
and bounded mutually disjoint sets $A_1, \dots, A_k \in \sB(\sX)$, we have  
\begin{align*}
P\left(\cap_{i=1}^k\set{N(A_i) = n_i}\right) = \prod_{i=1}^k\left(e^{-\Lambda(A_i)}\frac{\Lambda(A_i)^{n_i}}{n_i!}\right),\quad n \in \Z_+^k.
\end{align*}
\end{block} 
\begin{Definition}<2->
If the intensity measure $\Lambda$ of a Poisson process $S$ satisfies $\Lambda(A) = \lambda\abs{A}$ for all bounded $A \in \sB(\sX)$, then we call $S$ a \textbf{homogeneous Poisson point process} and $\lambda$ is its intensity.  
\end{Definition} 
\end{frame}

\begin{frame}{Equivalent characterizations}
\begin{block}{Theorem : Equivalences} 
\label{thm:equiv}
Following are equivalent for a simple counting process $N: \Omega \to {\Z_+}^{\sB(\sX)}$. 
\begin{enumerate}[i\_]
\item Process $N$ is Poisson with locally finite intensity measure $\Lambda$. 
\item For each bounded $A \in \sB(\sX)$, we have $P\set{N(A) = 0} = e^{-\Lambda(A)}$. 
\item %The intensity measure $\Lambda$ is bounded for bounded $A \in \sB$,  and 
For each bounded $A \in \sB(\sX)$,  the number of points $N(A)$ is a Poisson with parameter $\Lambda(A)$. 
\item Process $N$ has the completely independence property, and $\E N(A) = \Lambda(A)$ for all bounded sets $A \in \sB(\sX)$. 
\end{enumerate}
\end{block}
\end{frame}


\begin{frame}
\begin{block}{Proof}
We will show that $i\_ \implies ii\_ \implies iii\_ \implies iv\_ \implies i\_$. 
\begin{enumerate}
	\item[$i \implies ii\_$] It follows from the definition of Poisson point processes and definition of Poisson random variables.  
	\item[$ii \implies iii\_$] 
	From Corollary~\ref{cor:ExpPoisson}, we know that if void probabilities are exponential, then the marginal distributions are Poisson. 
	%From Theorem~\ref{thm:void} we know that void probabilities determine the entire distribution. 
\end{enumerate}
\end{block}
\end{frame}


\begin{frame}
\begin{block}{Proof continued}
\begin{enumerate}
\item[$iii \implies iv\_$] We will show this in two steps. 
\begin{itemize}
\item[Mean:] Since the distribution of random variable $N(A)$ is Poisson, it has mean $\E N(A) = \Lambda(A)$. 
\item[CIP: ] <2-> Consider a partition $A \in \sB^k$ for a bounded set $B\in\sB(\sX)$, then 
$\Lambda(B) = \Lambda(A_1) + \dots + \Lambda(A_k)$. 
Consider all partitions $n\in\Z_+^k$ of a non-negative integer $m \in \Z_+$, to write 
\EQ{
P\set{N(B)=m} = \sum_{n_1 + \dots + n_k = m}P\set{N(A_1) = n_1, \dots, N(A_k) = n_k}. 
}
Using the definition of Poisson distribution, we can write the LHS of the above equation as 
\EQ{
P\set{N(B) = m} 
= e^{-\Lambda(B)} \frac{\Lambda(B)^m}{m!} 
= \prod_{i=1}^ke^{-\Lambda(A_i)} \frac{(\sum_{i=1}^k\Lambda(A_i))^m}{m!}.
}
\end{itemize}
\end{enumerate}
\end{block}
\end{frame}


\begin{frame}
\begin{block}{Proof continued}
Since the expansion of $(a_1 + \dots + a_k)^m = \sum_{n_1 + \dots+n_k = m}\binom{m}{n_1, \dots, n_k}\prod_{i=1}^ka_i^{n_i}$, 
we get 
\begin{align*}
P\set{N(B) = m} 
&= \frac{1}{m!}\sum_{n_1 + \dots + n_ k =n}\binom{m}{n_1, \dots, n_k}\prod_{i=1}^ke^{-\Lambda(A_i)} \Lambda(A_i))^{n_i} \\
&= \sum_{n_1+\dots+n_k = n}\left(\prod_{i=1}^ke^{-\Lambda(A_i)} \frac{\Lambda(A_i))^{n_i}}{n_i!}\right). 
\end{align*}
Equating each term in the summation, we get $P\set{N(A_1) = n_1, \dots, N(A_k) = n_k} = \prod_{i=1}^kP\set{N(A_i) = n_i}$. 
\end{block}
\end{frame}


\begin{frame}
\begin{block}{Proof}
\begin{enumerate}
\item[$iv \implies i\_$] 
From Corollary~\ref{cor:ExpPoisson}, if the void probability is exponential with intensity measure $\Lambda$, then the marginal distribution if Poisson with the same intensity measure. 
We define $f:\sB(\sX)\to(-\infty, 0]$ by $f(A) \triangleq \ln P\set{N(A) = 0}$ for all bounded $A \in \sB(\sX)$. 
Then, we observe that for any partition $(A_1,\dots, A_k)$ of $A$, we have $f(\cup_{i=1}^kA_i) = \ln P\set{N(A)= 0} = \ln\prod_{i=1}^kP\set{N(A_i) = 0} = \sum_{i=1}^kf(A_i)$.
It follows that $-f:\sB(\sX)\to\R_+$ is an intensity measure, and $P\set{N(A)=0} = e^{f(A)}$. 
Since $\E N(A) = -f(A) = \Lambda(A)$, the result follows. 
\end{enumerate}
\end{block}
\end{frame}


\begin{frame}
\begin{cor}[Poisson process on the half-line]
A random process $N: \Omega \to \Z_+^{\R_+}$ indexed by time $t \in \Z_+$ is the counting process associated with a one-dimensional Poisson process $S: \Omega \to \R_+^\N$ having intensity measure $\Lambda$ iff
\begin{enumerate}[(a)]
\item<2-> Starting with $N_0 = 0$, the process $N_t$ takes a non-negative integer value for all $t \in \R_+$;
\item<3-> the increment $N_s - N_t$ is surely nonnegative for any $s \ge t$;
\item<4-> the increments $N_{t_1}, N_{t_2} - N_{t_1},  \dots , N_{t_n} - N_{t_{n-1}}$ are \textit{independent} for any 
$0 < t_1 < t_2 < \dots < t_{n-1} < t_n$; 
\item<5-> the increment $N_s - N_t$ is distributed as Poisson random variable with parameter $\Lambda(t, s]$ for $s \ge t$.  
\end{enumerate}
\end{cor}
\begin{itemize}
\item<6-> The Poisson process is homogeneous with intensity $\lambda$, iff in addition to conditions $(a), (b), (c)$, 
the distribution of the increment $N_{t+s} - N_t$ depends on the value $s \in \R_+$ but is independent of $t \in \R_+$. The increments are \textit{stationary}. 
\end{itemize}
\end{frame}

\begin{frame}
\begin{Proof} 
We have already seen that definition of Poisson processes implies all four conditions. 
Conditions $(a)$ and $(b)$ imply that $N$ is a simple counting process on the half-line, 
condition $(c)$ is the complete independence property of the point process, 
and condition $(d)$ provides the intensity measure. 
The result follows from the equivalence $iv\_$ in the Theorem on Equivalences. 
\end{Proof}
\end{frame}



  %P\set{X>s} = P(\set{ X > t+s}\given \set{ X>t}). 
 %\bar{F}(t+s) = \bar{F}(t)\bar{F}(s).
 %f(t+s) = f(t)f(s), \text{ for all } t,s \in \R_+
 %}
 %is $f(t) = e^{\theta t}$, where $\theta = \log f(1)$.
 %\end{xalignat*}

\end{document}
